% !TEX root = ../main.tex
\chapter{Introduction and Research Problem}
\label{ch:introduction}

Decentralized Finance (DeFi) has changed how capital is allocated on public blockchains: lending, borrowing and leveraged trading now run through smart contracts without traditional intermediaries \parencite{schar2021defi}. To use most of these services, an owner first authorizes a contract to move its tokens. The ERC-20 standard records each such authorization as an allowance \parencite{vogelsteller2015erc20}, and users, the software they sign with and other protocols all have to decide which spender contracts should receive one. Under pseudonymity that decision is hard to make well, and the central issue of this thesis is what the public record of earlier authorizations can say about a contract.

This thesis asks what ERC-20 allowance edges add to PageRank as a reputation score for smart contracts. The allowance edges come from the \texttt{Approval} events that \texttt{approve} and \texttt{permit} emit, and EndorseRank is PageRank computed on them with the edge weights of the transfer-based Adaptive Weighted PageRank (AWP) of \textcite{do2023awp}. Both are run on the same Arbitrum One contracts, beside a coupled walk over both edge types and the raw degrees that the walks smooth, and checked by the rank-correlation method used with AWP \parencite{do2023awp}, extended with paired contrasts, raw-degree baselines and two temporal holdouts, the second of them registered before its logs were extracted. The scores use on-chain data alone; Section~\ref{sec:research-problem} defines an on-chain reputation score together with the other main terms.

\dissertationSubheading[sec:background]{Background: DeFi, collateralization, and pseudonymity}

Public blockchains make programmable settlement possible without a central authority \parencite{nakamoto2008bitcoin}. A smart contract can encode the terms of a deal \parencite{szabo1997formalizing}, and Ethereum lets such a contract hold and move tokens under the rules of its code \parencite{buterin2014ethereum}. DeFi platforms built from these contracts clear trades on the blockchain instead of routing them through a custodian, and protocols stack them into decentralized markets for lending, automated market making and perpetual futures \parencite{werner2022sok}. On Arbitrum One and other L2 networks, fees are low enough that ordinary users transact often, and every transfer, approval and protocol event can be inspected address by address.

Even with this visibility, DeFi credit markets remain over-collateralized. Most protocols make a borrower or trader post more collateral than the position is worth, because the counterparty is a pseudonymous account whose owner has not been verified \parencite{harvey2021defi}. Pseudonymity is a basic property of public ledgers. Anyone can create addresses at almost no cost, and one entity can control many of them \parencite{douceur2002sybil}. Conventional credit ratings depend on personal histories that translate poorly to this setting \parencite{packin2024credit}, so protocols control risk through collateral reserves, liquidation mechanisms and protocol rules instead of personalized reputation. Researchers have turned to on-chain signals as a substitute for off-chain credit histories: transfer-graph centrality \parencite{do2023hodgerank}, network risk propagation \parencite{lin2025riskprop}, credit scoring based on prospect theory \parencite{hassija2020lending}, and machine learning pipelines over lending protocols \parencite{kotb2023credit}. All of these score borrowers or the accounts behind them.

Collateral protects a lending protocol against its borrowers, but it does nothing for an owner who must decide whether a contract may spend its tokens. In that decision the pseudonymous counterparty is the contract. Before a router, vault, aggregator or lending pool can move an owner's tokens, the owner grants it an allowance, which stays in force until it is spent, replaced or revoked \parencite{vogelsteller2015erc20}. Most approvals on Ethereum are for unlimited amounts \parencite{wang2022unlimited}, so an allowance granted to a malicious or later compromised contract exposes the owner's whole balance of that token. Phishing sites exploit this by leading users to sign transactions that hand their assets to attackers \parencite{he2023txphishscope}, and the toolkits that do so are now sold as a service \parencite{he2025drainer}. Protocols face the same decision when their own contracts approve other contracts, as a vault does when it lets a lending pool pull the tokens it supplies. A contract's code is public, but who deployed it, who can change it and how it will use an allowance often are not.

In practice, trust in a spender contract forms in a few ways. An owner usually decides at the moment its wallet software presents the approval for signature, and weaknesses in how wallets present and validate token approvals make that moment easy to exploit \parencite{adamczyk2025approval}. Code analysis judges a contract by its functions and can flag, for example, functions written for a rug pull \parencite{lin2024crpwarner}. Deployed reputation tools work at the level of persons. Human Passport collects weighted credentials, on-chain and off-chain, to decide whether an address is controlled by a distinct human, and it serves as a gate against Sybil attacks, not as a graded ordering \parencite{humanpassportnddocs}. That is a question about the holder of a key. It has no direct counterpart for a contract, whose behavior is fixed by its code and by whoever may upgrade it. None of these routes uses the public record of what other owners have decided: which contracts many owners have authorized, and how strongly those owners are endorsed in turn.

Graph ranking algorithms appear well suited to such a record. PageRank and its variants rank nodes by recursive endorsement, so that a node's score depends on the endorsements it receives from nodes that are endorsed in turn \parencite{page1999pagerank}. Web trust algorithms such as EigenTrust \parencite{kamvar2003eigentrust} and TrustRank \parencite{gyongyi2004trustrank} apply similar ideas to peer-to-peer trust and to web spam. \textcite{do2023hodgerank} use both PageRank and HodgeRank on a graph of transfers to judge how trustworthy Ethereum users are, and call the result a form of social credit. AWP \parencite{do2023awp} is another transfer-graph ranking algorithm. It weights each transfer by a bounded function of its value and a logistic function of its age, and it restarts the walk at addresses that sent recently; its authors check it by showing that it correlates with inbound transfer count and value, taken as proxies for economic importance. These studies make transfer-graph ranking the baseline against which a score built on another edge type has to be compared. They rank addresses by the payments they receive, however, and leave out explicit spending authorization, a separate on-chain action recorded as ERC-20 allowances \parencite{vogelsteller2015erc20}.

\dissertationSubheading[sec:motivation]{Motivation for on-chain reputation}

Four factors drive the need for a reputation score for spender contracts.
\begin{itemize}
\item Authorization decisions: owners, the software they sign with and the protocols that integrate other contracts all decide which spenders to approve. An ordinal signal of how other owners have treated a contract could inform these decisions, for instance as one input to the warning shown before an approval is signed, or to the list of contracts a protocol allows to integrate with it.
\item Auditability: curated lists and off-chain scores are opaque, but graph-based scores computed from public event data can be recomputed and audited by anyone, which answers part of the concern about black-box credit scoring in decentralized finance \parencite{packin2024credit}.
\item Operational cost: a score that is refreshed often, before each approval prompt for instance, should be cheap to recompute.
\item Standardization: other markets use screening ratios to tell counterparties apart, but DeFi has no standard ordinal language for doing so at the level of the contract. Chapter~\ref{ch:discussion} discusses where the scores studied here could serve that role and where they cannot.
\end{itemize}

Transfer-based AWP is public and auditable, but its graph grows with every payment \parencite{do2023awp}. Allowance graphs are a structurally different alternative. Whenever an owner calls \texttt{approve}, or signs a \texttt{permit} \parencite{lundfall2020permit}, the token contract records how much of that token a spender may transfer \parencite{vogelsteller2015erc20}, and each new approval supersedes the older ones, so the latest approval of each owner--spender pair and token stands for the history of that pair. Approvals are also rarer than transfers. On the same contracts the endorsement graph $G_{\text{approve}}$ therefore has far fewer edges than the transfer graph $G_{\text{transfer}}$, and the PageRank solve, whose cost per iteration grows linearly with $|E|$ \parencite{leskovec2020mining}, takes correspondingly less time (Section~\ref{sec:benchmark-results}). The same solver runs in both cases, so the difference comes from the sparser graph and is not an algorithmic improvement; Claim~C1 treats it as a property of the representation.

Efficiency aside, allowances also mean something different. A transfer might reflect routing, market making, or a passive receipt, whereas an allowance is an active grant of permission to spend funds. Treating allowances as endorsement edges carries PageRank's citation reading \parencite{page1999pagerank} over to a trust relationship that owners and protocols already depend on, and the endorsed node is the spender contract, the object that owners have to decide about. Whether this produces different rankings, and whether those rankings agree differently with authorization and transfer checks and with later activity, is examined on a matched cohort of $n=27{,}844$ contracts on Arbitrum One.\footnote{Matched means that every method scores the same contracts, a list fixed from the approval logs before any transfer was aggregated; Section~\ref{sec:data-sources} gives the rule.}

\dissertationSubheading[sec:research-problem]{Research problem}

AWP infers trust from transfer history through weighted sums and logistic time decay \parencite{do2023awp}. It gives each transfer an economic weight, but it (i) must aggregate the whole transfer history, (ii) models a high-value address as one that receives many inbound transfers, and (iii) assigns no role to spending authorization. EndorseRank, the score studied here, keeps AWP's edge weights and solver but builds the directed graph from the latest allowance of each owner--spender pair and token. The resulting graph is sparser, encodes a different signal and ranks a different side of each relationship. An allowance points from the owner to the spender, so EndorseRank scores the contracts that receive approvals, such as routers, vaults and lending pools, and each owner, an EOA or another contract, passes its endorsement to the spenders it approves.

The main research question is: what do ERC-20 allowance edges add to PageRank as a reputation score for smart contracts on Arbitrum One, compared with the transfer walk and with the raw degrees that both walks smooth? Research questions RQ1--RQ5 (Section~\ref{sec:research-questions}) break it down. Because both edge types live on the same addresses, a single walk can also mix them node by node with a fixed weight $\lambda$, which defines a Coupled PageRank (C-PR); at $\lambda=1$ and $\lambda=0$ it walks one layer alone. The mixture is worth testing because the two edge types record different acts on the same contracts. An operator that already runs a transfer walk could add allowance edges to it instead of running a second score, and the open question is whether the walk then carries the allowance signal without losing the transfer signal.

Agreement with statistics of a score's own edges, or prediction of more of the same edges, does not yet show that a score measures reputation in a substantive sense. A registered comparison, which adds no research question, therefore also uses an outcome built from neither edge type and observed after the freeze: the later liquidation avoidance of the GMX~V2 accounts that are contracts and trade in the label window (Section~\ref{sub:neutral-label-protocol}). Few such contracts traded in these months, so this comparison rests on small samples.

The scores appear in Table~\ref{tab:primary-methods-preview}. Table~\ref{tab:hrq-claim-map} lists hypotheses (H1--H3), research questions (RQ1--RQ5), and empirical claims (C1--C6) from Chapter~\ref{ch:results}.

\begin{table}[htbp]
\centering
\caption{The scores evaluated in this thesis (matched contract cohort, $n=27{,}844$).}
\label{tab:primary-methods-preview}
\fitwidth{%
\begin{tabular}{p{0.18\linewidth}p{0.28\linewidth}p{0.48\linewidth}}
\toprule
Score & Input graph & Role \\
\midrule
EndorseRank & Latest ERC-20 allowances ($G_{\text{approve}}$), weighted as in AWP; uniform restarts & Allowance delegation; the object of study in this research. A variant with AWP's restarts in proportion to activity is reported beside it \\
AWP & Transfers ($G_{\text{transfer}}$) & Transfer flow; the published transfer-graph method of \textcite{do2023awp} (Section~\ref{sub:awp-graph}) \\
C-PR & Both layers, raw amounts, mixed per node with weight $\lambda$ & Coupled operator used to test whether the two edge types combine (RQ5); at $\lambda=1$ and $\lambda=0$ it walks one layer alone. Its ablation is S-PR, a transfer walk restarted from the allowance walk \\
Raw degrees & $G_{\text{approve}}$ and $G_{\text{transfer}}$, counted per node & In-approve degree and transfer in-degree, the local counts that the PageRanks smooth: proxies in the same window and baselines out of window \\
\bottomrule
\end{tabular}
}
\end{table}

The main terms are defined as follows:

\begin{itemize}
\item \emph{DeFi (Decentralized Finance):} Financial services such as lending, trading and borrowing, implemented as smart contracts on public blockchains and generally run without a central operator in charge of custody or settlement \parencite{schar2021defi}.
\item \emph{Smart contract:} Self-executing program code that runs on a blockchain and carries out an agreement without a trusted intermediary \parencite{szabo1997formalizing}. In this thesis the allowance and transfer logs are emitted by token contracts, and the addresses that are scored are contracts as well.
\item \emph{Address, contract account and EOA:} An address identifies an account. A contract account is governed by the code it stores, whereas an externally owned account (EOA) is controlled by a private key \parencite{antonopoulos2018mastering}. An address counts as a contract here when \texttt{eth\_getCode} at one pinned block returns bytecode; no bytecode, or an EIP-7702 delegation designator \parencite{buterin2024eip7702}, makes it an EOA (Section~\ref{sub:cohorts}). EOAs enter the graphs as owners, senders and recipients and are never scored. No off-chain identity is assumed or recovered.
\item \emph{On-chain reputation score:} A number computed from on-chain data alone that expresses an address's relative position among the addresses it interacts with \parencite{do2023hodgerank}. In this thesis the scored addresses are contracts. Scores are ordinal and meant for ranking; their absolute scale has no meaning.
\item \emph{PageRank:} A graph ranking algorithm that spreads influence through a directed graph with a damping factor \parencite{page1999pagerank}. Adaptive Weighted PageRank (AWP) is PageRank on the transfer graph, with edges weighted by the value and age of each transfer and restarts at recently active senders \parencite{do2023awp}. EndorseRank applies AWP's edge weights to the latest-allowance graph and restarts uniformly (Section~\ref{sec:graph-construction}).
\item \emph{Coupled PageRank (C-PR):} One random walk over both the allowance and the transfer layer, which at each node mixes the two rows with weights $\lambda$ and $1-\lambda$ over the layers that node occupies (Section~\ref{sub:coupled-graph}). Endorsement-seeded PageRank (S-PR) keeps the transfer walk but restarts it from the allowance walk, as TrustRank restarts from a seed set \parencite{gyongyi2004trustrank}.
\item \emph{ERC-20 allowance / Approve / Permit:} Contract methods that record how many of the owner's tokens a spender may spend on the owner's behalf \parencite{vogelsteller2015erc20}. \texttt{Approve} is an ordinary on-chain call; EIP-2612 \texttt{permit} makes the same state change through a signature \parencite{lundfall2020permit}. For each owner--spender pair and token, EndorseRank uses the latest non-zero allowance.
\item \emph{Endorsement graph:} A directed graph whose edges denote an explicit delegation of spending authority, as opposed to an actual transfer of assets. An edge $u \to v$ means that the owner $u$, an EOA or a contract, has approved $v$ to spend tokens belonging to $u$ \parencite{vogelsteller2015erc20}. The spenders $v$ that are scored are contracts.
\item \emph{Domain alignment:} The degree to which a reputation ranking agrees with an observable proxy ranking, measured with Spearman's $\rho$ \parencite{spearman1904proof} and Kendall's $\tau$ \parencite{kendall1938rank}. Agreement with proxies built from a score's own edges is read as a check of the intended construct in the sense of \textcite{cronbach1955construct}. A method can agree more with the graph of its own construct than with a different check.
\end{itemize}

\dissertationSubheading[sec:supplementary-extension]{Temporal holdout}

The same-window allowance coefficients (Table~\ref{tab:alignment-hybrid}) are an intended-construct check: PageRank on $G_{\text{approve}}$ should agree with the in-approve degree of each spender, because it smooths that degree globally. To test whether a score conveys information beyond the same-window graph, the scores are frozen at $t_1$ (31~March~2026, 23:59:59~UTC) and compared with two labels counted from 1~April to 30~June~2026: the new owner--spender approval pairs that a spender contract gains (the endorsement construct) and its new transfer senders (the flow construct). This first window, W0, is evaluated on the spender cohort of every contract that holds at least one positive latest allowance at $t_1$ in the extracted approval data ($n=33{,}101$). Each PageRank is also compared with the raw degree it smooths at $t_1$, so the gain of the propagated score over the local count can be seen directly; these degree contrasts belong to no rule fixed in advance and were first computed after the W0 labels were known. W0 carries the primary part of rule A, the rule for the coupled operator fixed in the analysis plan, and an exploratory run on the GMX~V2 contract traders of the same months (Section~\ref{sec:holdout-protocol}).

A second window, W1, was registered at 16:49~UTC on 9~October~2026, after the W0 results were known and before any log dated after 30~June~2026 was extracted. The extraction script refuses to run unless the registration and the configuration are committed, and it records the hashes of both files. W1 freezes the scores at the end of the observation window (30~June~2026, 23:59:59~UTC), so they come from the same graphs as the same-window evaluation, and counts its labels from 1~July to 30~September~2026. It repeats the two spender labels on the contracts with a positive latest allowance at that freeze ($n=34{,}587$) and adds one for trader contracts: the share of closes that did not end in liquidation, for the GMX~V2 accounts that are contracts, appear in the freeze-date graph and close at least three positions in the window ($n=70$). The registration fixes rule B for the coupled operator, with its margin, three further contrasts reported outside the decision and two sensitivity analyses, and a comparison of allowance-side with transfer-side scores on that trader label, with the wording of each possible conclusion (Sections~\ref{sub:fresh-holdout} and~\ref{sub:neutral-label-protocol}). That comparison uses W0 as its secondary window, and the liquidation labels of the W0 traders had not been computed for any score when the registration was committed.

\begingroup
\setstretch{1}\small\setlength{\tabcolsep}{3.5pt}\captionsetup{font=normalsize}
\begin{longtable}{p{0.12\linewidth}p{0.46\linewidth}p{0.36\linewidth}}
\caption{Summary of hypotheses, research questions, and empirical claims.}
\label{tab:hrq-claim-map}\\
\toprule
Item & Addresses & Reported in \\
\midrule
\endfirsthead
\caption[]{(continued)}\\
\toprule
Item & Addresses & Reported in \\
\midrule
\endhead
\bottomrule
\endlastfoot
H1 / RQ1 & Do EndorseRank and AWP order the same contracts differently? & \S\ref{sec:rank-divergence}; inter-method $\tau$ with 95\% CI; Table~\ref{tab:tie-sensitivity}; Claim C3 \\
H2 / RQ2 & How closely does each score agree with the allowance checks and with the transfer checks? & Tables~\ref{tab:alignment-hybrid} and~\ref{tab:tie-sensitivity}; Claim C2 \\
H3 / RQ3 & What does each score cost to compute at equal $n$? & \S\ref{sec:benchmark-results}; Claim C1 \\
RQ4 & Do scores frozen at $t_1$ predict the later new approval pairs and new senders of spender contracts, and does any PageRank do better than the raw degree it smooths? & \S\ref{sec:holdout-results}; Tables~\ref{tab:holdout-spenders}, \ref{tab:holdout-diff}, \ref{tab:holdout-receiving}, \ref{tab:fresh-holdout} and~\ref{tab:fresh-posthoc}; Claim C4 \\
RQ5 & Does the coupled operator C-PR improve on a walk over the transfers alone, and does it rank contracts better than a walk over either layer alone? & \S\ref{sec:coupled-results}, \S\ref{sec:holdout-results} and~\S\ref{sub:fresh-results}; Tables~\ref{tab:alignment-hybrid}, \ref{tab:tau-diff-hybrid}, \ref{tab:holdout-diff}, \ref{tab:fresh-holdout} and~\ref{tab:fresh-contrasts}; Claim C5 \\
C1 & Solve time follows the edge count of each graph: EndorseRank's call takes less time than AWP's on its sparser graph, while C-PR and S-PR, which walk the transfer layer, take longer than AWP & \S\ref{sec:benchmark-results} \\
C2 & Same-window agreement: each score agrees most with the checks of the edge type it walks on, an intended-construct check that is partly built in; EndorseRank leaves almost no ties, the isolated-node rule changes no coefficient, and the C-PR mixtures lie between their two layers &Tables~\ref{tab:alignment-hybrid} and~\ref{tab:tie-sensitivity} \\
C3 & EndorseRank and AWP order the cohort differently & Table~\ref{tab:tie-sensitivity}; Figure~\ref{fig:rank-scatter} \\
C4 & Out of window, in contrasts computed after the labels were known, AWP ranks new transfer senders above the raw transfer in-degree in both windows, while EndorseRank as fixed in the plan ranks new approval pairs below the raw in-approve degree in both; its variant with activity restarts ranks them above the degree, a hypothesis for a registered test and not a finding &Tables~\ref{tab:holdout-spenders}, \ref{tab:holdout-diff}, \ref{tab:holdout-receiving} and~\ref{tab:fresh-posthoc} \\
C5 & One walk over both edge types does not beat its layers walked alone: rule A fails in both parts, and rule B is not met, because the gain over the transfer walk on new approval pairs replicates but non-inferiority fails on new senders and on the traders' liquidation-free close rate; against AWP it fails clearly; in the S-PR ablation, allowance edges help only when they carry the walk &Tables~\ref{tab:alignment-hybrid}, \ref{tab:tau-diff-hybrid}, \ref{tab:holdout-diff}, \ref{tab:fresh-holdout} and~\ref{tab:fresh-contrasts}; \S\ref{sub:fresh-results} \\
C6 & Registered comparison on an outcome built from neither edge type: no count-level lead of allowance over transfer on later liquidation avoidance in the primary window, the PageRank versions not resolved, and the approval count below the transfer count on the profitable and non-loss shares; small trader samples & Tables~\ref{tab:neutral-label-contrasts} and~\ref{tab:neutral-label-recipients} \\
\end{longtable}
\endgroup

\dissertationSubheading[sec:research-objectives]{Research objectives}

The research has five objectives, one for each of RQ1--RQ5. Each is stated with the measurement used to operationalize it and with the criterion that Chapter~\ref{ch:discussion} applies when judging whether the objective was met.

\begin{enumerate}
\item[O1.] Establish whether EndorseRank and AWP order the contracts of the same cohort differently (RQ1). Measurement: Kendall $\tau$ between the two score vectors on the matched cohort with a 95\% bootstrap interval. Criterion: the interval contains neither $0$ (unrelated orderings) nor $1$ (identical orderings).
\item[O2.] Assess how closely each score agrees with each construct in the same window (RQ2). Measurement: family-level Kendall $\tau$ with 95\% intervals for every score and construct on the matched cohort. Criterion: a score agrees with a construct when its interval on that family excludes $0$; agreement with the construct of the score's own edges is read as an intended-construct check, because those proxies come from the same events as the graph.
\item[O3.] Measure what each score costs to compute on the same contracts, and how runtime responds to the number of sampled contracts (RQ3). Measurement: solve time (mean and standard deviation of five timed runs after one warm-up run), peak memory, iterations, $|V|$ and $|E|$, all under one protocol. Criterion: the time ratio is reported only together with the ratio of $|E|$, so that any difference can be traced to the size of the graph.
\item[O4.] Test whether scores frozen before a label window predict the new approval pairs and new transfer senders that spender contracts gain in it, and whether any PageRank does better than the raw degree it smooths (RQ4). Measurement: Kendall $\tau$ with 95\% intervals between every score, the two raw degrees included, and both labels, on the spender contracts of two windows: W0, with scores frozen on 31~March~2026 and labels counted from April to June~2026, and W1, with scores frozen on 30~June~2026 and labels counted from July to September~2026 (Section~\ref{sec:supplementary-extension}); in each, paired contrasts of each PageRank against its own degree. Criterion: a score predicts a label when its interval excludes $0$, and a PageRank improves on its degree only when the interval of the paired difference lies above $0$. Neither the degree contrasts nor this criterion was fixed before the W0 labels were known.
\item[O5.] Test whether one walk across both edge types improves on a walk over the transfers alone, and whether it beats a walk over either layer alone (RQ5). Measurement: C-PR at $\lambda=0.5$ (with $0.25$ and $0.75$ as sensitivity values) and S-PR, compared through paired intervals on $\Delta\tau$ with the two layers walked alone and with EndorseRank and AWP, on the holdout and same-window labels. Criteria, each fixed before the data it is judged on: (a) rule A, fixed in the analysis plan of 9~October~2026 before any log was extracted, under which C-PR does no worse than the allowance layer walked alone on new approval pairs and better than the transfer layer walked alone on new senders in W0, and in the same window lies within the interval of the better layer on the transfer and allowance families and above both layers on Sybil stability; (b) rule B, registered later on 9~October~2026, after the W0 results were known and before the W1 logs were extracted, under which C-PR beats the transfer layer walked alone on new approval pairs in W1 and is not worse than it by more than a margin of $0.01$, set from the half-width of the corresponding W0 interval and rounded up, on new senders or on the trader contracts' liquidation-free close rate. A registered sensitivity analysis repeats (b) with AWP in place of the transfer layer; it does not change the rule. Every other contrast of C-PR and S-PR is exploratory: it was computed after the scores and labels it compares were known.
\end{enumerate}

\dissertationSubheading[sec:research-questions]{Research questions}

Three hypotheses go with the first three research questions. H1: on the same set of contracts, EndorseRank and AWP produce different orderings. H2: each score agrees with the checks built from the edges it walks on. H3: at the same value of $n$, a score's computation time follows the number of edges in its graph, so EndorseRank, on the sparser graph, is faster to compute than AWP. H1--H3 hold largely by construction and act as checks of the instruments. The falsifiable tests are RQ4, whose degree contrasts were first computed after the W0 labels were known and then repeated on the W1 labels outside the registered rules, and RQ5, answered by rules fixed before the data they are judged on (Section~\ref{sec:holdout-protocol}). Table~\ref{tab:hrq-claim-map} summarizes the questions.

\begin{enumerate}
\item[RQ1.] On the same set of Arbitrum One contracts, do EndorseRank and AWP produce different orderings?
\item[RQ2.] How closely does each score agree with current allowance checks and with current transfer checks (Spearman's $\rho$ and Kendall's $\tau$)?
\item[RQ3.] What does each score cost to compute on the same set of contracts, as measured by time, memory, and iterations?
\item[RQ4.] Given scores computed at a freeze date $t_1$, do the scores predict the new owner--spender approval pairs and the new transfer senders that spender contracts gain after $t_1$, and does any PageRank predict them better than the raw degree it smooths at $t_1$?
\item[RQ5.] Does a coupled PageRank running on both the allowance and the transfer layer improve on a walk over the transfers alone, such as AWP, and does it rank contracts better than a walk over either layer alone, out of window as well as in the same window?
\end{enumerate}

RQ4 uses later labels of the same two constructs because the contracts ranked here do not borrow, so no default or repayment label exists for them (Section~\ref{sub:non-claims}). RQ5 moves from comparing the two edge types to combining them. Rule A, which asks whether C-PR beats a walk over either layer alone, was fixed with $\lambda$ in the analysis plan and is judged on W0 and in the same window. The rule for improving on the transfer walk alone was written after the W0 results were known, so it is judged only on W1, whose logs were extracted after it was registered (Sections~\ref{sec:holdout-protocol} and~\ref{sub:fresh-holdout}).

\dissertationSubheading[sec:literature-summary]{Summary of the literature review}

Chapter~\ref{ch:literature} reviews four areas. In graph ranking, PageRank \parencite{page1999pagerank} ranks nodes by recursively propagated endorsement at a per-iteration cost linear in the number of edges \parencite{langville2006pagerank}, and EigenTrust \parencite{kamvar2003eigentrust} and TrustRank \parencite{gyongyi2004trustrank} adapt it to peer-to-peer trust and web spam with normalization and seed sets. In on-chain reputation, \textcite{do2023hodgerank} apply PageRank and HodgeRank to Ethereum transactions, and AWP \parencite{do2023awp} weights transaction edges by value and age and checks its ranking against inbound transfer count and value, statistics drawn from the same data as its graph. \textcite{lin2025riskprop} propagate account risk through the transaction network, and surveys of blockchain trust and reputation systems \parencite{bellini2020trust} and of smart-contract reputation designs \parencite{almasoud2020reputation} find explicit ratings and transfer volume to be the usual trust signals. Deployed DeFi tooling such as Human Passport instead aggregates on-chain and off-chain credentials to decide whether an address belongs to a distinct person \parencite{humanpassportnddocs}, and code analysis judges a contract by the functions it contains \parencite{lin2024crpwarner}. None of this work takes the ERC-20 allowance \parencite{vogelsteller2015erc20} as a ranking edge or ranks contracts by the approvals they receive. Address clustering uses approvals only to link addresses of one entity \parencite{victor2020clustering}, and an analysis of an approval-based fraud scheme shows that an approval can be obtained by deception \parencite{adamczyk2025approval}. In DeFi risk and measurement, systematizations classify risk by protocol layer \parencite{werner2022sok} and by the attacks seen in practice \parencite{zhou2023sok}, and wash trading \parencite{victor2021washtrading}, extractable value \parencite{qin2022bev}, scam tokens \parencite{xia2021scamtokens} and airdrop farming \parencite{luo2025airdrop} show that much transfer activity is artificial. In measurement theory, construct validity \parencite{cronbach1955construct} separates what a score measures from external criteria, and the bootstrap \parencite{efron1979bootstrap} gives distribution-free intervals for rank statistics. The gap (Section~\ref{sec:research-gap}) is an allowance-edge ranking of contracts evaluated on one cohort beside the transfer method and the raw degrees, with confidence intervals, paired differences, a fixed timing protocol and temporal holdouts, one of them registered before its data were extracted.

\dissertationSubheading[sec:methodology-summary]{Summary of the research methodology}

Chapter~\ref{ch:methodology} applies a quantitative, comparative design to public ledger data. An analysis plan, committed on 9~October~2026 before any log of the study was extracted, fixes the cohorts, scores, labels, statistics and rule A. ERC-20 \texttt{Approval} and \texttt{Transfer} logs on Arbitrum One from 1~October~2023 to 30~September~2026 are extracted from the public Google BigQuery dataset for the ego network of the matched cohort, whose rule and account-type check Section~\ref{sec:data-sources} defines. One power-iteration PageRank implementation scores the latest-allowance graph and the transfer graph, and the coupled operator runs on the same solver. Each score is compared with Spearman's $\rho$ and Kendall's $\tau$ against same-window transfer, allowance and Sybil-stability proxies and, out of window, against new approval pairs and new transfer senders on the spender contracts of W0 and W1, with the raw degrees at each freeze as baselines. Every coefficient has a 95\% interval from 400 paired contract resamples (seed 42), and method comparisons are paired. The registration of W1 adds rule B and a comparison of allowance-side with transfer-side scores on the later liquidation avoidance of GMX~V2 contract traders, a label built from neither edge type, in both windows (2{,}000 paired resamples; Bonferroni 97.5\% intervals for its two decision contrasts). Runtime, memory, iterations and graph size are recorded under one timing protocol. The code, configuration and machine-readable summaries behind every table are in a public repository (Section~\ref{sec:reproducibility}), and the study uses only public data (Section~\ref{sec:ethics}).

\dissertationSubheading[sec:contributions]{Contributions}

The thesis makes five contributions to the study of on-chain reputation and DeFi risk. EndorseRank and AWP both apply weighted PageRank \parencite{page1999pagerank} to a single edge type, each with its own restart distribution; only the coupled operator changes the transition rule. What is new is therefore an edge type and the objects it ranks, an evaluation design, an analytic account of what the allowance walk computes, a test of whether the two edge types combine, and a registered comparison on an outcome that neither edge type contains.
\begin{itemize}
\item Allowance edges as a ranking edge for contracts (EndorseRank): an endorsement graph in which an edge $u \to v$ stands for the latest in-window ERC-20 allowance, set by \texttt{approve} or \texttt{permit}, that the spender $v$ holds from the owner $u$, weighted as AWP weights a transfer. The edge records a delegation of spending rights, not a history of payments, and it points to the spender, so the score ranks the contracts that owners authorize. Approvals have been used to cluster the addresses of one entity \parencite{victor2020clustering}, but no reputation method reviewed ranks contracts on the approvals they receive (Section~\ref{sec:research-gap}).
\item Evaluation of every score against the counts it smooths, in the same window and in two holdout windows: every score is computed on one set of contracts and checked against identical transfer, allowance and Sybil-stability proxies under the protocol used with AWP \parencite{do2023awp}, with 95\% bootstrap intervals for every coefficient and paired intervals for every contrast \parencite{efron1979bootstrap}. Scores frozen before a label window are then tested against the new approval pairs and new transfer senders that spender contracts gain in it, beside the raw degrees that PageRank smooths, in W0 and again in W1, whose cohorts and labels were registered before its logs were extracted. The design separates same-window agreement from prediction and the propagated score from its local count. The degree contrasts belong to no fixed rule and were first computed after the W0 labels were known. On these data the two walks differ: the transfer walk ranked later senders above the raw transfer count in both windows, and the allowance walk as fixed in the plan did not rank later approvals above the raw approval count in either (Claim~C4).
\item Closed form of EndorseRank at depth one, and evidence of propagation: with uniform restarts and uniformly redistributed dangling mass, a spender whose owners receive no approvals themselves scores $c\,(1+d\,\delta(v))$, where $c$ is the common score of nodes without in-edges, $d$ the damping factor and $\delta(v)$ the owner-normalized in-strength, which divides each owner's endorsement over all the spenders it approves (Section~\ref{sub:theory-closed-form}). On the cohort's allowance graph the formula holds within 1\% for every depth-one spender but for only a minority of the spenders with a deeper in-neighbourhood, so multi-hop propagation is present in this graph. It did not, however, make EndorseRank rank later approvals better than the raw count.
\item Coupled operator and its tests: C-PR blends a raw-amount allowance layer and a raw-amount transfer layer node by node with a fixed weight $\lambda$, and S-PR is a seeded ablation. Rule A, fixed with $\lambda$ in the analysis plan before any log was extracted, asked C-PR to do no worse than the allowance layer and better than the transfer layer, each walked alone; rule B, registered before the logs of W1 were extracted, asked it to improve on the transfer layer walked alone. Both fail and are reported as failed, although the gain over the transfer walk on new approval pairs replicated in W1 (Claim~C5). The ablation, in contrasts that were not fixed in advance, shows that the allowance signal survives when allowance edges carry the walk and is lost when they only decide where it restarts.
\item Registered comparison on an outcome built from neither edge type: allowance-side and transfer-side scores are set against each other on the later liquidation avoidance of GMX~V2 contract traders, a label that enters neither graph, with the decision contrasts, their intervals and the wording of each conclusion registered before the July--September logs were extracted. In the primary window the registered conclusions are that there is no count-level lead of allowance over transfer on later liquidation avoidance and that the PageRank versions are not resolved; on the profitable and non-loss shares the approval count falls below the transfer count (Claim~C6). The trader samples are small, and the result is reported as registered.
\end{itemize}


\dissertationSubheading[sec:scope-limitations]{Scope and limitations}

The study covers Arbitrum One. Its observation window runs for 33 months, from 1~October~2023 to 30~June~2026, and the registered window W1 adds the quarter from 1~July to 30~September~2026. Its main sample is the $n=27{,}844$ contracts that received a non-zero approval from at least three distinct owners in the observation window, so the results apply to contracts of this kind on Arbitrum One in these months, and not to contracts approved by fewer owners. The holdout samples are the spender contracts with a positive latest allowance at each freeze ($n=33{,}101$ in W0 and $n=34{,}587$ in W1). Account type comes from \texttt{eth\_getCode} at one block read on 9~October~2026 (Section~\ref{sec:data-sources}), so an address is classified by its code at that block, not by its code at the time of each event. Only events with a cohort contract at one end are extracted, so each graph is the ego network of the cohort, and a path longer than one edge is seen only where every edge on it touches a cohort contract; multi-hop propagation is therefore measured on a truncated graph (Section~\ref{sub:theory-closed-form}). The comparison on outcomes built from neither edge type uses the GMX~V2 accounts that are contracts, and few of them traded in these months ($n=88$ in W0 and $n=70$ in W1), so every trader result has wide intervals (Section~\ref{sub:neutral-label-results}). Runtime scaling is measured on subsamples of the cohort and does not test graphs larger than the cohort's ego network.

The scores use no off-chain information: no audit reports, no curated lists of contracts and no identity credentials of the kind Human Passport collects. No score is tested as a detector of exploits, rug pulls or phishing, or against loan defaults (Section~\ref{sub:non-claims}), and no claim is made that the results generalize beyond Arbitrum One or the months studied. Amounts are raw token base units. EndorseRank and AWP pass them through a bounded function, so every approval or transfer of ten base units or more counts about the same, while in the raw layers of C-PR unlimited approvals decide the rows of many owners (Chapter~\ref{ch:methodology}). Adversarial manipulation of allowances or transfers is a real threat, because addresses are cheap to create \parencite{douceur2002sybil}. Chapter~\ref{ch:discussion} shows that a few Sybil owner addresses can lift a contract's EndorseRank above most of the cohort at little cost, while reaching the top of the cohort is expensive (Section~\ref{sec:sybil-cost-model}). The Sybil-adjusted proxies describe some manipulation patterns without protecting the score against them.

\dissertationSubheading[sec:roadmap]{Research roadmap}

Chapter~\ref{ch:literature} reviews graph-based on-chain reputation, ERC-20 allowance semantics and the reputation tooling deployed in DeFi, and develops the theoretical framework. Chapter~\ref{ch:methodology} defines the cohorts, the graphs and scores, the construct checks, the temporal holdouts W0 and W1, the timing protocol, and the standards for reproducibility and ethics. Chapter~\ref{ch:results} reports the benchmark (Claim~C1), the same-window alignment of every score (C2), rank divergence (C3), checks of damping, token set, sample definition and restarts, the temporal holdouts with the registered window (C4 and C5) and the registered comparison on outcomes built from neither edge type (C6); Appendix~\ref{ch:appendix-er-awp} collects the direct comparisons between EndorseRank and AWP. Chapter~\ref{ch:discussion} interprets the findings by research question and objective, works out their implications, models the cost of Sybil attacks on allowance graphs, and weighs threats to validity and limitations. Chapter~\ref{ch:conclusion} answers RQ1--RQ5, states the contribution to knowledge and sets out future work.
